\subsection{Материал Максвелла}

\inputfig{figures/fig-maxwell}

Пружина (элемент 1, модуль $\mu$) и демпфер (элемент 2, вязкость $\eta$)
соединены последовательно:
\[
  \Sdev = \Sdev_1 = \Sdev_2 ,
  \qquad
  \dev = \dev_1 + \dev_2 ,
\]
\[
  \Sdev_1 = 2\mu\,\dev_1 ,
  \qquad
  \Sdev_2 = 2\eta\,\dt{\dev}_2 .
\]
Дифференцируя первое по времени:
\[
  \dt{\dev} = \dt{\dev}_1 + \dt{\dev}_2
  = \frac{\dt{\Sdev}_1}{2\mu} + \frac{\Sdev_2}{2\eta}
  \quad\Longrightarrow\quad
  \boxed{\ \dt{\dev} = \frac{\dt{\Sdev}}{2\mu} + \frac{\Sdev}{2\eta}\ }
\]
— определяющее соотношение для материала Максвелла.

\subsubsection*{Ограничение модели}

Если задать $\Sdev$, то $\dev$ растёт линейно до $\infty$ (демпфер течёт
неограниченно). Поэтому задаём $\dev$, а $\Sdev$ ищем:
\[
  \dev = \dev_0 H(t)
  \quad \text{— кинематическое нагружение.}
\]
Вся деформация в начальный момент идёт в элемент 1 (пружину), так как иначе
скорость демпфера была бы бесконечной.

\subsubsection*{Релаксация}

\inputfig{figures/fig-relaxation-function}

Проинтегрируем строго. Сначала однородное уравнение ($\dt{\dev} = 0$):
\[
  \frac{\dt{\Sdev}}{2\mu} + \frac{\Sdev}{2\eta} = 0
  \quad\Longrightarrow\quad
  \tau\,\dt{\Sdev} + \Sdev = 0 ,
  \qquad
  \left[\ \tau = \frac{2\eta}{2\mu} = \frac{\eta}{\mu}
  \ \text{— время релаксации}\ \right]
\]
\[
  \Sdev = \tn C\, e^{-t/\tau} .
\]
Теперь неоднородное — метод вариации постоянной, $\Sdev = \tn C(t)e^{-t/\tau}$:
\[
  \frac{\dt{\tn C}(t)}{2\mu}e^{-t/\tau}
  \underbrace{{}- \tn C(t)e^{-t/\tau}\frac{1}{2\mu\tau}
  + \frac{\tn C(t)e^{-t/\tau}}{2\eta}}_{{}=\,0}
  = \dt{\dev}
\]
(средние слагаемые сокращаются, поскольку $2\mu\tau = 2\eta$), откуда
\[
  \dt{\tn C}(t) = 2\mu\, e^{t/\tau}\, \dt{\dev}(t)
  \quad\Longrightarrow\quad
  \tn C(t) = \int\limits_0^t 2\mu\, e^{\xi/\tau}\, \dt{\dev}(\xi)\,d\xi + \tn A ,
\]
где $\tn A = 0$ из нулевых начальных условий. Значит
\[
  \Sdev = \int\limits_0^t 2\mu\, e^{-(t-\xi)/\tau}\, \dt{\dev}(\xi)\,d\xi
        = \int\limits_0^t R(t-\xi)\,d\dev(\xi) ,
  \qquad
  \boxed{\ R(t) = 2\mu\, e^{-t/\tau}\ }
\]
— \textbf{функция релаксации}.

Пусть теперь
\[
  \dev = \dev_0 H(t) - \dev_0 H(t - t^*) .
\]
Напряжение скачком поднимается до $2\mu\dev_0$, релаксирует по экспоненте, а в
момент $t^*$ падает на ту же величину $2\mu\dev_0$ (скачок деформации вниз даёт
такой же по модулю скачок напряжения) и дальше возвращается к нулю снизу.

\inputfig{figures/fig-relax-pulse}

\subsection{Обобщённый материал Максвелла}

Параллельное соединение $N$ элементов Максвелла.

\inputfig{figures/fig-maxwell-generalized}

По аксиоме параллельного соединения девиаторы напряжений складываются при общей
деформации, поэтому складываются и функции релаксации:
\[
  R(t) = \sum_{k=1}^{N} 2\mu_k\, e^{-t/\tau_k} ,
  \qquad
  \tau_k = \frac{\eta_k}{\mu_k} .
\]

При $N = 2$ и $\tau_1 \ll \tau_2$ график падает от $2\mu_1 + 2\mu_2$ до уровня
$2\mu_2$ за время $\sim\tau_1$ (излом как в обобщённом материале
Кельвина–Фойгта), а затем медленно релаксирует к нулю за время $\sim\tau_2$.
